\section{Restricciones energéticas e iteración de hardware: el límite real de la competencia por desempeño}

\subsection{La restricción física del cómputo: al final es un presupuesto de energía}

Jensen señala explícitamente hacia la parte media y final: la escalabilidad del cómputo está limitada por la energía y el calor. Ya sea en el entrenamiento o en la inference, en esencia siempre hay que equilibrar el consumo de energía, la disipación térmica, la latencia y el rendimiento.

\begin{knowledgebox}{Por qué la «mejora del desempeño» depende cada vez más del diseño de sistemas}
Aumentar únicamente la frecuencia de un solo chip ya no basta para sostener las cargas de trabajo de IA; es necesario depender de:
\begin{itemize}
    \item la optimización de la arquitectura del chip (el grado de paralelismo, la jerarquía de memoria)
    \item la interconexión de alta velocidad (la comunicación entre GPU o nodos)
    \item la programación de software a nivel de sistema (compiladores, runtime, estrategias de paralelización)
\end{itemize}
Por eso la competencia actual es una «competencia de ingeniería de sistemas», no una «competencia de un solo dispositivo».
\end{knowledgebox}

\subsection{Del centro de datos al dispositivo personal: el descenso de la forma de la AI supercomputer}

En la entrevista, Jensen mostró sistemas de etapas tempranas y mencionó la evolución de equipos de investigación de alto costo hacia dispositivos personales de menor umbral de acceso. Esto refleja la misma tendencia: el cómputo de IA descenderá hacia un grupo más amplio de desarrolladores y estudiantes, tal como ocurrió en su momento con la computación personal.

\begin{figure}[H]
\centering
\includegraphics[width=0.88\textwidth]{frame-dgx-a.jpg}
\caption{En la entrevista se recuerda una escena temprana de entrega de una AI supercomputer \protect\footnotemark}
\end{figure}
\footnotetext{Intervalo de tiempo en el video: 00:36:58--00:37:06.}

\begin{figure}[H]
\centering
\includegraphics[width=0.88\textwidth]{frame-geforce-a.jpg}
\caption{El presentador muestra una forma de AI supercomputer personal desplegable en escritorio \protect\footnotemark}
\end{figure}
\footnotetext{Intervalo de tiempo en el video: 00:52:33--00:52:45.}

\begin{importantbox}{Conclusión clave}
La «AI democratization» no es solo la apertura de modelos, sino también la migración de la forma del cómputo y de la cadena de herramientas de desarrollo hacia un rango accesible para las personas. Por eso la velocidad de la innovación educativa e industrial experimentará una segunda aceleración.
\end{importantbox}

\subsection{Resumen de la sección}
La competencia en IA ha pasado de «quién tiene el modelo más grande» a «quién puede seguir ofreciendo, bajo restricciones energéticas, una mayor eficiencia de sistema y un umbral de uso más bajo».

